\chapter{Complexity — Poly Time Model}
%%% generated by the blueprint pipeline from TCSlib.Complexity.Randomized.PolyTimeModel
\section{Overview}

Declarations of the Lean module \texttt{TCSlib.Complexity.Randomized.PolyTimeModel}.

\section{Declarations}

\begin{definition}[Polynomial-time verifier model]
\lean{Randomized.polyTimeModel}
\statementsource{Vad12}{
  pdf-bf2ce5422088-p020-b004,
  pdf-bf2ce5422088-p021-b003,
  pdf-bf2ce5422088-p021-b005,
  pdf-bf2ce5422088-p023-b001
}
\statementsource{HILL99}{
  pdf-130039751e9f-p009-b003,
  pdf-130039751e9f-p009-b004
}
\statementsource{arora-barak-ch07-randomized-computation}{pdf-26d95080dbde-p002-b003}
\label{Randomized.polyTimeModel}
\leanok
\uses{Complexity.P, Randomized.VerifierModel, Turing.pairEncode}
A verifier model consists of a notion of efficiency for option-valued two-input
verifiers and a notion of efficiency for Boolean-valued three-input predicates. In the
polynomial-time model, an option-valued verifier is efficient exactly when there are
languages $V_1,V_0$ in $P$ such that, for every input $x$ and string $r$, it returns
true exactly when the pair encoding of $(x,r)$ belongs to $V_1$, and returns false
exactly when that encoding belongs to $V_0$; a three-input predicate is efficient
exactly when there is a language $V$ in $P$ such that its value on $(x,u,v)$ is true
exactly when the nested pair encoding of $(x,u,v)$ belongs to $V$.
\end{definition}

\begin{definition}[Taking a bounded pair-encoded slice]
\lean{Randomized.sliceTake}
\label{Randomized.sliceTake}
\leanok
\uses{Complexity.pairFstD, Complexity.pairSndD, Randomized.polyLen, Turing.pairEncode}
For natural numbers $a$ and $k$ and a Boolean list $z$, write $(u,v)$ for the pair
decoded from $z$. The slice is the pair encoding of $u$ together with the prefix of $v$
of length $	ext{polyLen}(a,k,|u|)$.
\end{definition}

\begin{definition}[Slice drop]
\lean{Randomized.sliceDrop}
\label{Randomized.sliceDrop}
\leanok
\uses{Complexity.pairFstD, Complexity.pairSndD, Randomized.polyLen, Turing.pairEncode}
For natural numbers $a$ and $k$ and a finite Boolean list $z$, decode $z$ into a pair of
Boolean lists $(u,v)$. The slice-drop list is the pair encoding of $u$ and the list
obtained from $v$ by deleting its first $	ext{polyLen}(a,k,|u|)$ entries.
\end{definition}

\begin{definition}[Decoded prefix extraction]
\lean{Randomized.takePrefixByLen}
\label{Randomized.takePrefixByLen}
\leanok
\uses{Complexity.pairFstD, Complexity.pairSndD}
Given a Boolean list $p$, decode it into two Boolean lists $u$ and $v$ using the pair
decoder. The defined list is the prefix of $v$ of length equal to the length of $u$.
\end{definition}

\begin{definition}[Dropping a decoded prefix]
\lean{Randomized.dropPrefixByLen}
\label{Randomized.dropPrefixByLen}
\leanok
\uses{Complexity.pairFstD, Complexity.pairSndD}
For a Boolean list $p$, let $u$ and $v$ be its first and second decoded components,
respectively. The definition assigns the list obtained from $v$ by dropping its first
$|u|$ entries.
\end{definition}

\begin{theorem}[Polynomial-time length-indexed prefix extraction]
\lean{Randomized.polyTimeComputable_takePrefixByLen}
\label{Randomized.polyTimeComputable_takePrefixByLen}
\leanok
\uses{Complexity.PolyTimeComputable, Complexity.PrefixByLength.take, Complexity.polyTimeComputable_takePrefixByLength, Randomized.takePrefixByLen}
\difficulty{1}
The function sending a bit string $p$ to the prefix of its decoded second component
whose length is the length of its decoded first component is polynomial-time computable;
failed decoding yields empty components.
\end{theorem}

\begin{theorem}[Polynomial-Time Prefix Removal]
\lean{Randomized.polyTimeComputable_dropPrefixByLen}
\label{Randomized.polyTimeComputable_dropPrefixByLen}
\leanok
\uses{Complexity.PolyTimeComputable, Complexity.PrefixByLength.drop, Complexity.polyTimeComputable_dropPrefixByLength, Randomized.dropPrefixByLen}
\difficulty{1}
For a bit string $z$, decode its pair, take the second component (or the empty string if
decoding fails), and remove from its beginning as many bits as the length of the first
component (or zero if decoding fails). The resulting function on bit strings is
computable by a deterministic finite-state Turing machine in time $C(n+1)^c$ on every
input of length $n$, for some natural numbers $C$ and $c$.
\end{theorem}

\begin{theorem}[Polynomial-time prefix truncation]
\lean{Randomized.polyTimeComputable_sliceTake}
\label{Randomized.polyTimeComputable_sliceTake}
\leanok
\uses{Complexity.PolyTimeComputable, Complexity.PolyTimeComputable.comp, Complexity.PolyTimeComputable.pairEncode, Complexity.pairFstD, Complexity.pairFstD_pairEncode, Complexity.pairSndD, Complexity.pairSndD_pairEncode, Complexity.polyTimeComputable_pairFstD, Complexity.polyTimeComputable_pairSndD, Complexity.polyTimeComputable_polyUnary, Randomized.polyLen, Randomized.polyTimeComputable_takePrefixByLen, Randomized.sliceTake, Randomized.takePrefixByLen, Turing.pairEncode}
\difficulty{4}
For all natural numbers $a$ and $k$, the following function on bit strings is
polynomial-time computable: decode the input into components $x$ and $y$, using empty
components when decoding fails, and return the encoding of $x$ paired with the prefix of
$y$ of length $a(|x|+1)^k$.
\end{theorem}

\begin{theorem}[Polynomial-time suffix truncation]
\lean{Randomized.polyTimeComputable_sliceDrop}
\label{Randomized.polyTimeComputable_sliceDrop}
\leanok
\uses{Complexity.PolyTimeComputable, Complexity.PolyTimeComputable.comp, Complexity.PolyTimeComputable.pairEncode, Complexity.pairFstD, Complexity.pairFstD_pairEncode, Complexity.pairSndD, Complexity.pairSndD_pairEncode, Complexity.polyTimeComputable_pairFstD, Complexity.polyTimeComputable_pairSndD, Complexity.polyTimeComputable_polyUnary, Randomized.dropPrefixByLen, Randomized.polyLen, Randomized.polyTimeComputable_dropPrefixByLen, Randomized.sliceDrop, Turing.pairEncode}
\difficulty{4}
For all natural numbers $a$ and $k$, the following function on bit strings is
polynomial-time computable: decode the input into components $x$ and $y$, using empty
components when decoding fails, and return the encoding of $x$ paired with the suffix of
$y$ obtained by dropping its first $a(|x|+1)^k$ bits.
\end{theorem}

\begin{theorem}[Closure under sequential race]
\lean{Randomized.polyTimeModel_closedUnderRace}
\label{Randomized.polyTimeModel_closedUnderRace}
\leanok
\uses{Complexity.compl_mem_P, Complexity.inter_mem_P, Complexity.pairFstD_pairEncode, Complexity.pairSndD_pairEncode, Complexity.preimage_mem_P, Randomized.ClosedUnderRace, Randomized.boolVerifier, Randomized.polyLen, Randomized.polyTimeComputable_sliceDrop, Randomized.polyTimeComputable_sliceTake, Randomized.polyTimeModel, Randomized.raceVerifier, Randomized.sliceDrop, Randomized.sliceTake, Turing.pairEncode}
\difficulty{5}
For every pair of Boolean predicates $M_1,M_2$ and natural numbers $a,k$, if the
verifiers associated to both predicates are efficient in the model, then the verifier
that applies $M_1$ to the first $a(|x|+1)^k$ bits of $r$, returns true if that accepts,
and otherwise applies $M_2$ to the remaining suffix and returns false if that accepts
(returning no answer if it does not) is efficient in the model.
\end{theorem}

\begin{theorem}[Closure under fixed answers]
\lean{Randomized.polyTimeModel_closedUnderAnswerIs}
\label{Randomized.polyTimeModel_closedUnderAnswerIs}
\leanok
\uses{Complexity.compl_mem_P, Randomized.ClosedUnderAnswerIs, Randomized.polyTimeModel}
\difficulty{4}
For every verifier $M$ and Boolean value $b$, if $M$ is efficient in the polynomial-time
verifier model, then the verifier that accepts $(x,r)$ exactly when $M(x,r)=b$ is also
efficient in that model.
\end{theorem}

\begin{theorem}[Closure under majority repetition]
\lean{Randomized.polyTimeModel_closedUnderMajority}
\label{Randomized.polyTimeModel_closedUnderMajority}
\leanok
\uses{Complexity.blockAt, Complexity.compl_mem_P, Complexity.mem_P_of_blockMajority, Complexity.pairFstD, Complexity.pairFstD_pairEncode, Complexity.pairSndD_pairEncode, Randomized.ClosedUnderMajority, Randomized.boolVerifier, Randomized.majorityVerifier, Randomized.polyLen, Randomized.polyTimeModel, Turing.MultiTapeTM.indicator, Turing.pairEncode}
\difficulty{5}
For every Boolean predicate $M$ and natural numbers $a,k,a',k'$, if the verifier that
returns $M(x,r)$ is efficient in the model, then so is the verifier that returns true
when more than half of the first $a'(|x|+1)^{k'}$ consecutive blocks of $r$, each of
length $a(|x|+1)^k$, make $M$ return true.
\end{theorem}

\begin{theorem}[Closure under existential repetition]
\lean{Randomized.polyTimeModel_closedUnderAny}
\label{Randomized.polyTimeModel_closedUnderAny}
\leanok
\uses{Complexity.blockAt, Complexity.compl_mem_P, Complexity.mem_P_of_blockAny, Complexity.pairFstD, Complexity.pairFstD_pairEncode, Complexity.pairSndD_pairEncode, Randomized.ClosedUnderAny, Randomized.anyVerifier, Randomized.boolVerifier, Randomized.polyLen, Randomized.polyTimeModel, Turing.pairEncode}
\difficulty{4}
For every Boolean predicate $M$ and natural numbers $a,k,a',k'$, if the verifier that
returns $M(x,r)$ is efficient in the model, then so is the verifier that returns true
when at least one of the first $a'(|x|+1)^{k'}$ consecutive blocks of $r$, each of
length $a(|x|+1)^k$, makes $M$ return true.
\end{theorem}

\begin{theorem}[Closure under complementation]
\lean{Randomized.polyTimeModel_closedUnderNot}
\label{Randomized.polyTimeModel_closedUnderNot}
\leanok
\uses{Randomized.ClosedUnderNot, Randomized.boolVerifier, Randomized.polyTimeModel}
\difficulty{3}
For every Boolean predicate $M$, if the verifier that returns $M(x,r)$ is efficient in
the model, then so is the verifier that returns the Boolean complement of $M(x,r)$.
\end{theorem}

\begin{theorem}[Closure under shifted disjunction]
\lean{Randomized.polyTimeModel_closedUnderShiftOr}
\label{Randomized.polyTimeModel_closedUnderShiftOr}
\leanok
\uses{Complexity.blockAt, Complexity.mem_P_of_blockXorAny, Complexity.pairFstD, Complexity.pairFstD_pairEncode, Complexity.pairSndD, Complexity.pairSndD_pairEncode, Randomized.ClosedUnderShiftOr, Randomized.boolVerifier, Randomized.polyLen, Randomized.polyTimeModel, Randomized.shiftOrVerifier, Turing.pairEncode}
\difficulty{4}
For every Boolean predicate $M$ and natural numbers $a,k,a',k'$, if the verifier
associated to $M$ is efficient in the model, then the predicate on three bit strings
that returns true when, for some $i<a'(|x|+1)^{k'}$, applying $M$ to $x$ and the
coordinatewise exclusive-or of $v$ with the block of $u$ beginning at position
$i\,a(|x|+1)^k$ and having length $a(|x|+1)^k$ returns true is an efficient two-witness
predicate in the model. The exclusive-or is taken pairwise over the two lists up to the
length of the shorter list.
\end{theorem}

\begin{theorem}[Circuits for polynomial-time verifiers]
\lean{Randomized.polyTimeModel_verifierHasCircuits}
\label{Randomized.polyTimeModel_verifierHasCircuits}
\leanok
\uses{BoolCircuit.DAGCircuit, BoolCircuit.DAGCircuit.IsFaninTwo, BoolCircuit.DAGCircuit.eval, BoolCircuit.DAGCircuit.pairEncode, BoolCircuit.DAGCircuit.pairEncode_eval, BoolCircuit.DAGCircuit.pairEncode_isFaninTwo, BoolCircuit.DAGCircuit.pairEncode_size, BoolCircuit.DAGCircuit.size, BoolCircuit.DAGCircuitFamily.pairEncode_eval_eq_true_iff, BoolCircuit.pairEncodeInput, Complexity.P_subset_PPoly, Language.InPPoly, Randomized.VerifierHasCircuits, Randomized.boolVerifier, Randomized.polyLen, Randomized.polyTimeModel, Turing.pairEncode}
\difficulty{6}
For every Boolean predicate $M$ and natural numbers $a,k$, if the verifier associated to
$M$ is efficient in the polynomial-time verifier model, then there exist natural numbers
$A,K$ such that, for every $n$ and every bit string $r$ of length $a(n+1)^k$, there is a
well-formed Boolean circuit with $n$ input bits, at most two inputs to each gate, and
size at most $A(n+1)^K$, whose output on each $v\in\{0,1\}^n$ equals $M(v,r)$. Here
well-formed means that no gate lists an input more than once and every negation gate has
exactly one input.
\end{theorem}

\begin{theorem}[Polynomial-time randomized to polynomial-size circuits]
\lean{Randomized.adleman_polyTime}
\label{Randomized.adleman_polyTime}
\leanok
\uses{Language.InPPoly, Randomized.InBPP, Randomized.adleman, Randomized.polyTimeModel, Randomized.polyTimeModel_closedUnderMajority, Randomized.polyTimeModel_verifierHasCircuits}
\difficulty{1}
Let $L$ be a language of bit strings. If $L$ has a polynomial-time randomized verifier
whose acceptance probability is at least $2/3$ on members and whose rejection
probability is at least $2/3$ on nonmembers, then $L$ is decidable by a family of
fan-in-two circuits of polynomial size.
\end{theorem}

\begin{theorem}[Second-level polynomial hierarchy characterization]
\lean{Randomized.inSigma2_polyTimeModel_iff}
\label{Randomized.inSigma2_polyTimeModel_iff}
\leanok
\uses{Complexity.PolyHierarchy.UnaryPT, Complexity.PolyHierarchy.mem_altClass_of_normal, Complexity.PolyHierarchy.unaryPT_poly, Complexity.SigmaP, Complexity.altQuant, Complexity.mem_SigmaP_two_iff_exists_forall, Complexity.polyTimeComputable_id, Complexity.qStep, Randomized.InSigma2, Randomized.polyLen, Randomized.polyTimeModel, Turing.pairEncode}
\difficulty{5}
For every language $L$ of bit strings, membership in the class defined by an efficient
two-witness predicate with polynomially bounded witness lengths is equivalent to
membership in $\Sigma_2^P$. Explicitly, there are a polynomial-time language $V$ and
natural numbers $C,c$ such that, for every bit string $x$, $x\in L$ if and only if there
is a bit string $u$ of length $C(|x|+1)^c$ for which every bit string $v$ of length
$C(|x|+1)^c$ satisfies $\langle\langle x,u\rangle,v\rangle\in V$.
\end{theorem}

\begin{theorem}[Characterization of zero-error randomized polynomial time]
\lean{Randomized.zpp_eq_rp_inter_corp_polyTime}
\label{Randomized.zpp_eq_rp_inter_corp_polyTime}
\leanok
\uses{Randomized.InCoRP, Randomized.InRP, Randomized.InZPP, Randomized.inZPP_iff_inRP_and_inCoRP, Randomized.polyTimeModel, Randomized.polyTimeModel_closedUnderAnswerIs, Randomized.polyTimeModel_closedUnderAny, Randomized.polyTimeModel_closedUnderRace}
\difficulty{1}
For every language $L$ of Boolean strings, $L$ belongs to zero-error randomized
polynomial time if and only if it belongs to both randomized polynomial time and
co-randomized polynomial time, with all three classes defined using the polynomial-time
verifier model.
\end{theorem}

\begin{theorem}[Polynomial-time BPP in the second level]
\lean{Randomized.sipser_gacs_polyTime}
\label{Randomized.sipser_gacs_polyTime}
\leanok
\uses{Complexity.PiP, Complexity.SigmaP, Randomized.InBPP, Randomized.inSigma2_polyTimeModel_iff, Randomized.polyTimeModel, Randomized.polyTimeModel_closedUnderMajority, Randomized.polyTimeModel_closedUnderNot, Randomized.polyTimeModel_closedUnderShiftOr, Randomized.sipser_gacs}
\difficulty{1}
If a language $L$ has a polynomial-time randomized verifier that uses polynomially many
random bits and answers correctly with probability at least $2/3$ on every input, then
$L$ belongs to both $\Sigma_2^P$ and $\Pi_2^P$. Here these classes are defined by
polynomially bounded alternating quantifiers over witnesses and a polynomial-time
decidable predicate.
\end{theorem}
